\documentclass{article}
\usepackage{geometry}
 \geometry{
 a4paper,
 left=20mm,
 right=20mm,
 top=20mm,
 }
\usepackage[affil-it]{authblk}
\usepackage{hyperref}
% \renewcommand{\familydefault}{\sfdefault}
\usepackage{newtxtext,newtxmath}
\usepackage{graphicx, cancel, booktabs, tabularx, amsmath, amsfonts, amssymb, bm, placeins, float,mathrsfs,url}
\usepackage{float}
\usepackage[utf8]{inputenc}
\renewcommand{\thefigure}{R\arabic{figure}}
\renewcommand{\thetable}{R\arabic{table}}
\usepackage{enumitem}   % custom enumeration
\usepackage{hyperref}   % autoref
\usepackage{booktabs}
% \usepackage{svg}
\usepackage{caption}
\usepackage{subcaption}
\usepackage{graphicx}
\usepackage{longtable}
\usepackage{adjustbox}

\hypersetup{
colorlinks = true,
linkcolor = [rgb]{0.87,0.18,0.22},
citecolor = [rgb]{0.0,0.60,0.33},
urlcolor = [rgb]{0.20, 0.30, 0.54}}
% \definecolor{eggplant}{RGB}{180,33,147}

\renewcommand{\theequation}{Revision Eq.~\arabic{equation}}

% \newcounter{responsecounter}
% \setcounter{responsecounter}{-1} 

% \newenvironment{response}[1][]{
%   \refstepcounter{responsecounter}
%   \par\medskip
%   \noindent
%   \textbf{Response~\theresponsecounter.}~#1\par
%   \smallskip
% }{\par\medskip}
\usepackage{orcidlink}
\usepackage{xcolor} % defines colors 
\usepackage[most]{tcolorbox} %  colored and framed text boxes
\definecolor{block-gray}{gray}{0.95}
\newtcolorbox{xresponse}{%
	empty,
    borderline west = {4pt}{0pt}{gray},
    boxrule = 0pt,
    boxsep = 0pt,
    breakable,
    colback = block-gray,
    enhanced,
    frame hidden,
    left skip = 0pt,
    notitle,
    parbox = false,
    sharp corners,
}
\renewcommand\Authfont{\small}
\renewcommand\Affilfont{\small\itshape}

% set a counter for the environment so that \autoref can be used
\newcounter{responseCounter}
\newenvironment{response}[1][]{\begin{xresponse}\refstepcounter{responseCounter}\textbf{Response~\theresponseCounter. #1}}{\end{xresponse}}

% set the \autoref abbreviation for the response environment
\def\responseCounterautorefname{Resp.}
\def\responseCounterautorefnameplural{Resps.}

% use \Autoref{one,two,three} for multiple references in one command
% from https://tex.stackexchange.com/a/183682
\makeatletter
% define a macro \Autoref to allow multiple references to be passed to \autoref
\newcommand\Autoref[1]{\@first@ref#1,@}
\def\@throw@dot#1.#2@{#1}% discard everything after the dot
\def\@set@refname#1{%    % set \@refname to autoefname+s using \getrefbykeydefault
    \edef\@tmp{\getrefbykeydefault{#1}{anchor}{}}%
    \xdef\@tmp{\expandafter\@throw@dot\@tmp.@}%
    \ltx@IfUndefined{\@tmp autorefnameplural}%
         {\def\@refname{\@nameuse{\@tmp autorefname}s}}%
         {\def\@refname{\@nameuse{\@tmp autorefnameplural}}}%
}
\def\@first@ref#1,#2{%
  \ifx#2@\autoref{#1}\let\@nextref\@gobble% only one ref, revert to normal \autoref
  \else%
    \@set@refname{#1}%  set \@refname to autoref name
    \@refname~\ref{#1}% add autoefname and first reference
    \let\@nextref\@next@ref% push processing to \@next@ref
  \fi%
  \@nextref#2%
}
\def\@next@ref#1,#2{%
   \ifx#2@ and~\ref{#1}\let\@nextref\@gobble% at end: print and+\ref and stop
   \else, \ref{#1}% print  ,+\ref and continue
   \fi%
   \@nextref#2%
}
\makeatother

% ################################
%          Title and author
% ################################

% ======= Title =======
\title{\large \textbf{{\sc Response Letter}\\Title of your article}}

% ======= Authors =======
\author[1]{AAA~\orcidlink{1111-1111-1111-1111}}
\author[1,2]{AAA\thanks{\href{mailto:xxx@whu.edu.cn}{xxx@whu.edu.cn}}~\orcidlink{0000-0000-0000-0000}}
% ======= Affiliations =======
\affil[1]{School of Integrated Circuits, Wuhan University, Wuhan 430072, China}
\affil[2]{Key Laboratory of Artificial Micro- and Nano-Structures of Ministry of Education, School of Physics and Technology, Wuhan University, Wuhan 430072, China}


% ======= Date =======
\date{\small\today}
% ################################
%          Begin document
% ################################

\begin{document}
\maketitle

% ################################
%         Intro/Letter
% ################################
% - Direct response, revision/rebuttal letter
\begin{description}
\item[Submission ID:] XXX
\item[Journal:] \textit{Physical Review Letters}
\item[Title:] ``\textit{XXX}''
\end{description}

\noindent
\textbf{Dear Editor(s) \& Referees},

We sincerely thank the editor(s) for forwarding the referees' reports on our manuscript. We greatly appreciate the referees' constructive comments, which have provided valuable insights to enhance our manuscript. In the following, we address the referees' comments in detail and outline our responses and the corresponding revisions made in the manuscript. We hope that our revised manuscript meets the standards for publication.
If the referees have any further questions, please feel free to contact us at the address below.

Best regards,
\begin{flushright}
	\textbf{Prof. XXX}, on behalf of the co-authors.\\
	\textit{School of XXX, Wuhan University}
\end{flushright}
\clearpage
% ################################
%    Own summary of reviews
% ################################
% - Main issues and responses
% - Summaries of each individual review
% \input{sections/owns_all}
% \newpage

% ################################
%    Full reviews and responses
% ################################

\tableofcontents
\clearpage

% Insert all reviews that you received
% Then, interleave the reviews for your responses

\section*{Comments \& Responses} \label{sec:reviewsFull}

% \input{sections/review_sum}

% ################################
%          Referee \#A
% ################################
% Insert the first review that you received

% \newcounter{responsecounter}
% \setcounter{responsecounter}{-1}

% \newenvironment{response}[1][]{
%   \refstepcounter{responsecounter}
%   \par\medskip
%   \noindent
%   \textbf{Response~\theresponsecounter.}~#1\par
%   \smallskip
% }{\par\medskip}

\section{Referee \#A}
\label{sec:rev:1}

\subsection{Summary}
XXX

\textbf{Reply: }We thank the referee for their recognition and encouraging comments.

\subsection{Referee \#A Comment I}

XXX

\begin{response}
\label{res:rev1:com1}

We thank the referee for pointing out this potential issue. 

\end{response}

\clearpage
% ################################
%          Referee \#B
% ################################
% Insert the second review that you received

% \newcounter{responsecounter}
% \setcounter{responsecounter}{0}

\section{Referee \#B}
\label{sec:rev:2}

\subsection{Summary}
xxx

\subsection{Referee \#B Comment I}

\textbf{Inadequate benchmarking against state-of-the-art models.}

xxx

\begin{response} 
\label{res:rev2:com1}

We thank the referee for the constructive suggestion~\cite{cgcnn}.

\end{response}



\newpage

% ################################
%        Difference view
% ################################
% Difference view between previous and current submissions, see:
% https://github.com/ftilmann/latexdiff/
% https://www.overleaf.com/learn/latex/Articles/Using_Latexdiff_For_Marking_Changes_To_Tex_Documents
% https://ctan.org/pkg/latexdiff

% NOTE: uncomment \usepackage{pdfpages} in the Packages section above
%\includepdf[pages=-]{diff_versions/diff_final.pdf}

\bibliographystyle{naturemag}  % or plainnat, abbrvnat, etc.
\bibliography{ref}

\end{document}
